\subsection{bin-fold 尾和阈值阶梯、window-\texorpdfstring{$6$}{6} 有限容量与 Fourier--R\'enyi 审计}\label{subsec:conclusion-foldbin-tail-capacity-fourier-renyi-audit}
沿用定义 \ref{def:K-of-m}、命题 \ref{prop:fold-bin-fiber-tail}、推论 \ref{cor:fold-bin-m6-fiber-hist}、定理 \ref{thm:fold-collision-spectrum} 与推论 \ref{cor:fold-collision-scale-diverges} 的记号。由这些内部结果可进一步把 tail-sum budget、window-\(6\) 的有限窗可逆容量、以及折叠分布的 Fourier--entropy 审计压缩为一组完全显式的终端后果。其共同结构是：bin-fold 纤维在固定末位层上由 admissible tail sums 的阈值截断精确控制；在 \(m=6\) 时，这一离散阶梯塌缩为闭式容量曲线与三能级热力学；而在渐近尺度上，碰撞缺口又强制真正参与的非平凡 Fourier 模数无界增长，从而排除 \(M^{-1/2}\) 自然尺度上的总变差收敛。

\begin{theorem}[bin-fold 纤维的尾和阈值阶梯律]\label{thm:conclusion-foldbin-tail-sum-threshold-staircase}
对任意 \(m\ge 1\) 与 \(w=(w_1,\dots,w_m)\in X_m\)，记
\[
V_m(w):=\sum_{k=1}^{m}w_kF_{k+1},
\qquad
b:=w_m,
\]
并令 \(K(m)\) 满足
\[
F_{K(m)+1}\le 2^m-1<F_{K(m)+2}.
\]
定义 admissible tail-sum 集合
\[
\mathcal T_{m,b}
:=
\left\{
\sum_{k=m+1}^{K(m)}c_kF_{k+1}:\ 
c_k\in\{0,1\},\ 
c_kc_{k+1}=0,\ 
bc_{m+1}=0
\right\}.
\]
则
\[
\boxed{\;
d_m^{\mathrm{bin}}(w)
=
\#\bigl\{t\in \mathcal T_{m,b}:\
t\le 2^m-1-V_m(w)\bigr\}.\;}
\]
因此，在固定末位 \(b\) 时，\(d_m^{\mathrm{bin}}(w)\) 只依赖于 deficit
\[
\Delta_m(w):=2^m-1-V_m(w),
\]
并且作为 \(V_m(w)\) 的函数是单调非增的阶梯函数；等价地，其阈值位置恰由
\[
V=2^m-1-t,
\qquad t\in\mathcal T_{m,b},
\]
完全决定。
\end{theorem}
\begin{proof}
由命题 \ref{prop:fold-bin-fiber-tail}，\(\Fold_m^{\mathrm{bin}}{}^{-1}(w)\) 与满足下列三类约束的尾部位串 \((c_{m+1},\dots,c_{K(m)})\) 一一对应：相邻约束 \(c_kc_{k+1}=0\)、边界约束 \(bc_{m+1}=0\)、以及阈值约束
\[
V_m(w)+\sum_{k=m+1}^{K(m)}c_kF_{k+1}\le 2^m-1.
\]
对任一 admissible 尾串定义其权重和
\[
t(\mathbf c):=\sum_{k=m+1}^{K(m)}c_kF_{k+1}.
\]
由于 admissible 二进制串给出的 Fibonacci 展开满足 Zeckendorf 型唯一性，映射 \(\mathbf c\mapsto t(\mathbf c)\) 在 admissible 尾串集合上是单射，其像正是 \(\mathcal T_{m,b}\)。因此纤维多重度恰为满足
\[
t(\mathbf c)\le 2^m-1-V_m(w)=\Delta_m(w)
\]
的 tail-sum 个数，也即
\[
d_m^{\mathrm{bin}}(w)
=
\#\bigl\{t\in \mathcal T_{m,b}:t\le \Delta_m(w)\bigr\}.
\]
这就给出了所述计数公式。

在固定 \(b\) 时，\(\mathcal T_{m,b}\) 与 \(m\) 固定后不再依赖于 \(w\)，故 \(d_m^{\mathrm{bin}}(w)\) 只依赖于 \(\Delta_m(w)\)。而函数
\[
\Delta\longmapsto \#\bigl\{t\in \mathcal T_{m,b}:t\le \Delta\bigr\}
\]
显然关于 \(\Delta\) 单调非减，于是关于 \(V=2^m-1-\Delta\) 单调非增，并且只有当 \(\Delta\) 穿过某个 \(t\in\mathcal T_{m,b}\) 时才会改变取值；等价地，其阈值位置正是 \(V=2^m-1-t\)。证毕。
\end{proof}

\begin{corollary}[window-\texorpdfstring{$6$}{6} 的连续容量曲线完全闭式]\label{cor:conclusion-window6-continuous-capacity-piecewise-closed}
取 \(m=6\)，并定义连续容量
\[
\mathcal C_6^{\mathrm{cont}}(T):=\sum_{w\in X_6}\min\bigl(d_6^{\mathrm{bin}}(w),T\bigr)
\qquad(T\ge 0).
\]
则有完全闭式
\[
\boxed{\;
\mathcal C_6^{\mathrm{cont}}(T)=
\begin{cases}
21T, & 0\le T\le 2,\\[2pt]
16+13T, & 2\le T\le 3,\\[2pt]
28+9T, & 3\le T\le 4,\\[2pt]
64, & T\ge 4.
\end{cases}\;}
\]
若进一步记最优 \(B\)-bit 预言机成功率
\[
\mathrm{Succ}_6(B):=\frac{\mathcal C_6^{\mathrm{cont}}(2^B)}{64},
\]
则
\[
\boxed{\;
\mathrm{Succ}_6(0)=\frac{21}{64},
\qquad
\mathrm{Succ}_6(1)=\frac{21}{32},
\qquad
\mathrm{Succ}_6(B)=1\ \ (B\ge 2).\;}
\]
\end{corollary}
\begin{proof}
由推论 \ref{cor:fold-bin-m6-fiber-hist}，window-\(6\) 的纤维直方图为
\[
2:8,\qquad 3:4,\qquad 4:9.
\]
因此直接有
\[
\mathcal C_6^{\mathrm{cont}}(T)
=
8\min(2,T)+4\min(3,T)+9\min(4,T).
\]
逐段化简即得所示分段闭式。再取 \(T=2^B\) 并除以总微态数 \(2^6=64\)，便得到三条成功率公式。证毕。
\end{proof}

\begin{theorem}[非平凡活跃 Fourier 模数必然无界增长]\label{thm:conclusion-foldbin-active-fourier-modes-diverge}
记
\[
M:=F_{m+2},
\qquad
\mathsf{Col}_m:=\sum_{r\in \ZZ/(M\ZZ)}p_m(r)^2,
\qquad
p_m(r):=\frac{d_m(r)}{2^m},
\]
\[
\mathcal Z_m:=\{k\in \ZZ/(M\ZZ):\ \widehat d_m(k)=0\},
\qquad
A_m:=M-1-|\mathcal Z_m|.
\]
则有下界
\[
\boxed{\;
A_m\ge M\Bigl(\mathsf{Col}_m-\frac1M\Bigr)
=M\,\mathsf{Col}_m-1.\;}
\]
特别地，由推论 \ref{cor:fold-collision-scale-diverges} 的 \(M\,\mathsf{Col}_m\to+\infty\)，可得
\[
\boxed{\;
A_m\longrightarrow +\infty.\;}
\]
\end{theorem}
\begin{proof}
记归一化 Fourier 幅度
\[
a_m(k):=\frac{|\widehat d_m(k)|}{2^m}.
\]
由定理 \ref{thm:fold-collision-spectrum}，
\[
\mathsf{Col}_m-\frac1M
=
\frac1M\sum_{k\ne 0}a_m(k)^2.
\]
另一方面，由三角不等式
\[
|\widehat d_m(k)|
\le \sum_r d_m(r)
=2^m,
\]
故对每个 \(k\) 都有 \(0\le a_m(k)\le 1\)。因此
\[
\sum_{k\ne 0}a_m(k)^2
\le
\#\{k\ne 0:\ a_m(k)\ne 0\}
=
A_m.
\]
两式合并便得
\[
M\Bigl(\mathsf{Col}_m-\frac1M\Bigr)\le A_m.
\]
这就证明了第一条。

再由推论 \ref{cor:fold-collision-scale-diverges}，
\[
M\,\mathsf{Col}_m\longrightarrow +\infty,
\]
故右侧下界 \(M\,\mathsf{Col}_m-1\) 也趋于 \(+\infty\)，从而 \(A_m\to+\infty\)。证毕。
\end{proof}

\begin{corollary}[自然尺度 \texorpdfstring{$M^{-1/2}$}{M^-1/2} 上的总变差收敛不可能成立]\label{cor:conclusion-foldbin-natural-scale-tv-impossible}
沿用定理 \ref{thm:conclusion-foldbin-active-fourier-modes-diverge} 的记号，并令均匀分布
\[
u(r):=\frac1M
\qquad
\bigl(r\in \ZZ/(M\ZZ)\bigr).
\]
则有
\[
\boxed{\;
\sqrt{M}\,\|p_m-u\|_2
=
\sqrt{M\Bigl(\mathsf{Col}_m-\frac1M\Bigr)}
=
\sqrt{M\,\mathsf{Col}_m-1}
\longrightarrow +\infty.\;}
\]
从而总变差距离
\[
\mathrm{TV}(p_m,u):=\frac12\|p_m-u\|_1
\]
满足
\[
\boxed{\;
\sqrt{M}\,\mathrm{TV}(p_m,u)\longrightarrow +\infty.\;}
\]
因而，在自然尺度 \(M^{-1/2}\) 上不可能有 \(p_m\) 向均匀分布 \(u\) 的总变差收敛。
\end{corollary}
\begin{proof}
直接展开可得
\[
\|p_m-u\|_2^2
=
\sum_r p_m(r)^2-\frac{2}{M}\sum_r p_m(r)+\sum_r\frac1{M^2}
=
\mathsf{Col}_m-\frac1M,
\]
因为 \(\sum_r p_m(r)=1\) 且 \(M\) 个点上 \(\sum_r M^{-2}=1/M\)。故
\[
\sqrt{M}\,\|p_m-u\|_2
=
\sqrt{M\Bigl(\mathsf{Col}_m-\frac1M\Bigr)}
=
\sqrt{M\,\mathsf{Col}_m-1}.
\]
再由推论 \ref{cor:fold-collision-scale-diverges} 知 \(M\,\mathsf{Col}_m\to+\infty\)，于是上式趋于 \(+\infty\)。

最后，由 \(\|x\|_1\ge \|x\|_2\)，得到
\[
\sqrt{M}\,\mathrm{TV}(p_m,u)
=
\frac{\sqrt{M}}{2}\|p_m-u\|_1
\ge
\frac{\sqrt{M}}{2}\|p_m-u\|_2,
\]
其右端已经趋于 \(+\infty\)，故结论成立。证毕。
\end{proof}

\begin{theorem}[window-\texorpdfstring{$6$}{6} 的全 R\'enyi 热力学完全闭式]\label{thm:conclusion-window6-full-renyi-thermodynamics-closed}
令 \(\mu_6\) 为 \(m=6\) 的归一化折叠输出分布：
\[
\mu_6(w):=\frac{d_6^{\mathrm{bin}}(w)}{64}
\qquad (w\in X_6).
\]
则对任意 \(q>0\) 且 \(q\neq 1\)，R\'enyi 熵具有精确闭式
\[
\boxed{\;
H_q(\mu_6)
:=
\frac{1}{1-q}\log\!\left(\sum_{w\in X_6}\mu_6(w)^q\right)
=
\frac{1}{1-q}
\left[
\log\!\bigl(8\cdot 2^q+4\cdot 3^q+9\cdot 4^q\bigr)-q\log 64
\right].\;}
\]
相应的 Shannon 熵为
\[
\boxed{\;
H(\mu_6)
:=
-\sum_{w\in X_6}\mu_6(w)\log\mu_6(w)
=
6\log 2
-\frac{1}{64}\Bigl(8\cdot 2\log 2+4\cdot 3\log 3+9\cdot 4\log 4\Bigr).\;}
\]
若进一步令 escort 分布
\[
\pi_{6,q}(w):=\frac{\mu_6(w)^q}{\sum_{u\in X_6}\mu_6(u)^q},
\]
并把它按纤维大小推前到 \(\{2,3,4\}\)，则
\[
\boxed{\;
\pi_{6,q}(d=2)=\frac{8\cdot 2^q}{8\cdot 2^q+4\cdot 3^q+9\cdot 4^q},\;}
\]
\[
\boxed{\;
\pi_{6,q}(d=3)=\frac{4\cdot 3^q}{8\cdot 2^q+4\cdot 3^q+9\cdot 4^q},\;}
\]
\[
\boxed{\;
\pi_{6,q}(d=4)=\frac{9\cdot 4^q}{8\cdot 2^q+4\cdot 3^q+9\cdot 4^q}.\;}
\]
\end{theorem}
\begin{proof}
由推论 \ref{cor:fold-bin-m6-fiber-hist}，
\[
d_6^{\mathrm{bin}}(w)=2\ \text{出现 }8\text{ 次},\qquad
3\ \text{出现 }4\text{ 次},\qquad
4\ \text{出现 }9\text{ 次}.
\]
因此
\[
\sum_{w\in X_6}\mu_6(w)^q
=
8\left(\frac{2}{64}\right)^q
+4\left(\frac{3}{64}\right)^q
+9\left(\frac{4}{64}\right)^q
=
\frac{8\cdot 2^q+4\cdot 3^q+9\cdot 4^q}{64^q}.
\]
代入 R\'enyi 熵定义便得到第一条公式。

Shannon 熵则直接计算为
\[
\begin{aligned}
H(\mu_6)
&=
-\sum_{w}\frac{d_6^{\mathrm{bin}}(w)}{64}
\log\!\left(\frac{d_6^{\mathrm{bin}}(w)}{64}\right)\\
&=
\log 64-\frac{1}{64}\sum_{w}d_6^{\mathrm{bin}}(w)\log d_6^{\mathrm{bin}}(w),
\end{aligned}
\]
而直方图给出
\[
\sum_{w}d_6^{\mathrm{bin}}(w)\log d_6^{\mathrm{bin}}(w)
=
8\cdot 2\log 2+4\cdot 3\log 3+9\cdot 4\log 4,
\]
故
\[
H(\mu_6)
=
6\log 2-\frac{1}{64}\Bigl(8\cdot 2\log 2+4\cdot 3\log 3+9\cdot 4\log 4\Bigr).
\]

最后，escort 分布满足
\[
\pi_{6,q}(w)\propto \mu_6(w)^q\propto d_6^{\mathrm{bin}}(w)^q.
\]
因此把 \(\pi_{6,q}\) 按纤维大小聚合后，只需对三个能级分别按 \(8\cdot 2^q\)、\(4\cdot 3^q\)、\(9\cdot 4^q\) 归一化，即得到后三条闭式。证毕。
\end{proof}

\noindent
以上结果把 bin-fold 口径下看似分散的四类对象压缩为同一条可审计链：尾和 admissibility 产生严格阈值阶梯，window-\(6\) 的有限窗容量因此塌缩为一条三段线性曲线并给出最优位预算的闭式基准，碰撞尺度的无界增长又迫使非平凡 Fourier 活跃模数无限增殖并排除 \(M^{-1/2}\) 自然尺度上的总变差均匀化，而同一三能级直方图则把全部 R\'enyi 熵、Shannon 熵与 escort 质量函数完全闭式化。由此，fold multiplicity 不再只是单一的组合计数对象，而同时表现为阈值账本、容量相图、频谱活化律与有限热力学单元。

\endinput
